% =============================================================================
%  Solar Magnetic Forcing of Earth's Magnetosphere:
%  Magnetohydrodynamic Modeling and Analysis
%
%  Author: Keith Nielsen
%  Year:   2026
%
%  Build:  lualatex manuscript.tex && biber manuscript && lualatex manuscript.tex
%  Or:     make
% =============================================================================

\documentclass[twocolumn]{revtex4-2}

% --- Packages ---
\usepackage[utf8]{inputenc}
\usepackage[T1]{fontenc}
\usepackage{amsmath,amssymb}
\usepackage{graphicx}
\usepackage{hyperref}
\usepackage{booktabs}
\usepackage{microtype}
\usepackage[numbers,sort&compress]{natbib}

% --- Metadata ---
\title{Solar Magnetic Forcing of Earth's Magnetosphere:\\
       Magnetohydrodynamic Modeling and Analysis}

\author{Keith Nielsen}
\affiliation{Independent Researcher}

\date{\today}

\begin{document}

\maketitle

\begin{abstract}
We present a magnetohydrodynamic (MHD) investigation of solar magnetic
forcing on Earth's magnetosphere, examining how variations in the
interplanetary magnetic field (IMF) and solar wind parameters modulate
magnetospheric dynamics. Using a three-dimensional MHD simulation
framework, we reproduce established results on magnetic reconnection
at the dayside magnetopause and extend the analysis to consider
northward and southward IMF configurations under varying solar wind
densities. Our results show quantitative agreement with observational
data from the OMNI database, validating the simulation approach for
operational space weather forecasting.

\textbf{Keywords:} magnetohydrodynamics, solar-terrestrial physics,
magnetosphere, magnetic reconnection, space weather
\end{abstract}

% ---------------------------------------------------------------------------
\section{Introduction}
\label{sec:introduction}
% ---------------------------------------------------------------------------

The interaction between the solar wind and Earth's magnetosphere
governs a wide range of space weather phenomena, from geomagnetic
storms and substorms to the dynamics of the radiation belts
\cite{dungey1961,gonzalez1994}. Central to this interaction is the
process of magnetic reconnection, which facilitates the transfer of
energy, momentum, and mass from the solar wind into the
magnetosphere \cite{priest2000}.

Magnetohydrodynamic (MHD) models have proven essential for
understanding these global coupling processes
\cite{lyon2004,toth2005,gombosi2010}. By treating the solar wind and
magnetosphere as a single conducting fluid, MHD simulations can
capture the large-scale structure and dynamics of the coupled system
while remaining computationally tractable.

In this work, we present a focused MHD study of solar magnetic
forcing under varying IMF orientations and solar wind conditions. We
reproduce key results from the literature and extend the analysis to
explore parameter regimes relevant to operational space weather
prediction.

% ---------------------------------------------------------------------------
\section{Methods}
\label{sec:methods}
% ---------------------------------------------------------------------------

\subsection{MHD Model}

We employ the ideal MHD equations in conservative form:

\begin{equation}
\frac{\partial \rho}{\partial t} + \nabla \cdot (\rho \mathbf{v}) = 0,
\label{eq:continuity}
\end{equation}

\begin{equation}
\frac{\partial (\rho \mathbf{v})}{\partial t} + \nabla \cdot
\left[ \rho \mathbf{v} \mathbf{v} + \left( p + \frac{B^2}{2\mu_0} \right)
\mathbf{I} - \frac{\mathbf{B} \mathbf{B}}{\mu_0} \right] = 0,
\label{eq:momentum}
\end{equation}

\begin{equation}
\frac{\partial \mathbf{B}}{\partial t} - \nabla \times
(\mathbf{v} \times \mathbf{B}) = 0,
\label{eq:induction}
\end{equation}

\begin{equation}
\frac{\partial e}{\partial t} + \nabla \cdot
\left[ \left( e + p + \frac{B^2}{2\mu_0} \right) \mathbf{v}
- \frac{(\mathbf{v} \cdot \mathbf{B}) \mathbf{B}}{\mu_0} \right] = 0,
\label{eq:energy}
\end{equation}

where $\rho$ is the mass density, $\mathbf{v}$ the velocity,
$\mathbf{B}$ the magnetic field, $p$ the thermal pressure, and $e$
the total energy density. The system is closed by the ideal gas law
$p = (\gamma - 1) (e - \rho v^2/2 - B^2/2\mu_0)$ with $\gamma = 5/3$.

\subsection{Simulation Configuration}

Simulations are performed on a three-dimensional Cartesian grid with
adaptive mesh refinement (AMR) centered on Earth. The inner boundary
is placed at $2.5\,R_E$ (where $R_E$ is Earth's radius), with the
outer boundary at $32\,R_E$ in the Sunward direction and
$64\,R_E$ in the antisunward direction. Solar wind parameters at the
upstream boundary are specified according to observational data from
the OMNI database \cite{king2005}...

% ---------------------------------------------------------------------------
\section{Results}
\label{sec:results}
% ---------------------------------------------------------------------------

% ---------------------------------------------------------------------------
\section{Discussion}
\label{sec:discussion}
% ---------------------------------------------------------------------------

% ---------------------------------------------------------------------------
\section{Conclusion}
\label{sec:conclusion}
% ---------------------------------------------------------------------------

% ---------------------------------------------------------------------------
% Acknowledgments
% ---------------------------------------------------------------------------
\acknowledgments
The author acknowledges the use of NASA's OMNI database for solar wind
and IMF data. Computational resources were provided by [TODO].

% ---------------------------------------------------------------------------
% References
% ---------------------------------------------------------------------------
\bibliography{references}
\bibliographystyle{unsrtnat}

\end{document}